\section{规则三的协整配对筛选算法}\label{sec:pairalgo}

本节给出规则三配对筛选的完整算法步骤。筛选分七步，每步都给出判据、实现要点与拒绝率实测。实现在 \texttt{code/rule3\_pairs.py}，可直接运行（\texttt{python rule3\_pairs.py --screen}）。

\subsection{算法的整体结构}

筛选的目的是从产业链内的两两组合中，找出「价差确实平稳、且回复速度落在可交易区间」的那些对。七步的顺序不可调换：先分组降维、再初筛降噪、再严格检验、最后估计参数。

\begin{quote}\ttfamily\small
S1 产业链分组 $\to$ 生成候选对\\
S2 相关性初筛（去伪配对）\\
S3 Engle--Granger 协整检验（OLS + ADF）\\
S4 OU 参数估计（AR(1) $\to$ $\theta$、半衰期）\\
S5 半衰期区间筛选 $[5,60]$ 交易日\\
S6 滚动重估（形成期 60 日 / 交易期 21 日）\\
S7 生成可交易对表（对冲比率、均值、标准差、半衰期）
\end{quote}

\subsection{S1：产业链分组与候选对生成}

\textbf{判据}：只在同一产业链内配对。跨产业链的配对数不足信——两条序列之间缺少经济意义上的长期均衡关系，协整容易被伪回归误判。

\textbf{实现}：按表~\ref{tab:chains}~的五条产业链展开。设第 $c$ 条链有 $n_c$ 个品种，则该链产生 $\binom{n_c}{2}$ 对；去重后汇总。

\textbf{实测}：五条产业链共 118 个候选对。

\subsection{S2：相关性初筛}

\textbf{判据}：形成期两序列对数价格的相关性 $\rho\ge0.6$。

\textbf{理由}：这一步不是统计检验，而是计算量与稳健性的过滤。相关性过低的配对在后续 ADF 检验中几乎必然被拒，先剔除以节省计算；更重要的是，低相关配对即便偶然通过协整检验，其价差波动也极大，交易风险不可控。

\textbf{实现}：
\begin{equation}
\rho_{AB}=\mathrm{Corr}\bigl(\ln P^A_t,\ \ln P^B_t\bigr),\qquad t\in\text{形成期}.
\end{equation}

\textbf{实测}：16992 次检验中被此步拒绝 8760 次（$51.6\%$），是最主要的过滤环节。

\subsection{S3：Engle--Granger 协整检验}

\textbf{判据}：对残差做 ADF 检验，$p<0.05$。这是本算法唯一的「资格检验」，其依据是反转类策略的条件性特征（论文对偶七）：\emph{必须先证明价差会回归，再交易}。

\textbf{实现}（两步法\cite{englegranger1987}）：

\begin{enumerate}[leftmargin=2.2em]
  \item 对形成期做 OLS：$\ln P^A_t=\alpha+\beta\ln P^B_t+e_t$，取 $\hat\beta$；
  \item 对残差 $e_t$ 做 ADF 检验（含常数项、滞后阶数由 AIC 自动选择\cite{dickeyfuller1979}），
        统计量需小于 $5\%$ 临界值；等价地 $p<0.05$。
\end{enumerate}

\textbf{实现要点}：OLS 用最小二乘直接求解（\texttt{lstsq}），ADF 用统计库的 \texttt{adfuller}。残差必须做去均值处理，否则 ADF 的常数项设定与协整回归的截距重复。

\textbf{已知局限}：Engle--Granger 两步法在样本量较小时检验功效偏低（论文附录 A.1 讨论的 $\hat\phi$ 下偏问题同源）。若形成期只有 60 个观测，ADF 的临界值并非标准分布，严格做法应使用 Engle--Granger 专用的响应面临界值（比标准 ADF 临界值更负）。本实现在此采用标准 ADF 临界值，属\emph{偏宽松}的近似，会略微高估通过率。

\textbf{实测}：被此步拒绝 4942 次（$29.1\%$）。通过者的平均 ADF $p$ 值为 $0.0203$。

\subsection{S4：OU 参数估计}

\textbf{判据}：残差可拟合为稳定 AR(1)，即 $0<\phi<1$。

\textbf{实现}：对残差序列拟合
\begin{equation}
e_{t+1}=c+\phi e_t+u_t ,
\end{equation}
由 OLS 得到 $\hat\phi$（同样可用极大似然，二者在正态假设下渐近等价）。再换算为连续时间参数：
\begin{equation}
\theta=-\frac{\ln\hat\phi}{\Delta},\qquad
\mu=\frac{c}{1-\phi},\qquad
H=\frac{\ln 2}{\theta}\quad\text{（半衰期）}.
\end{equation}

\textbf{为何必须检查 $0<\phi<1$}：$\phi\ge1$ 意味着残差非平稳（价差不回归），$\phi\le0$ 意味着残差是高频振荡而非缓慢回归，两者都使半衰期失去意义。论文附录指出 $\hat\phi$ 在小样本下有向下偏误，因此短形成期估出的半衰期偏短——这一点在 S5 的区间下界设置中有直接后果（见下）。

\textbf{实测}：仅 1 次被此步拒绝，说明在通过 ADF 的基础上 $\hat\phi$ 几乎总落在 $(0,1)$ 内。

\subsection{S5：半衰期区间筛选}

\textbf{判据}：$H\in[5,\,60]$ 交易日。

\textbf{两侧边界的含义}：
\begin{itemize}[leftmargin=2.2em]
  \item \textbf{下界 $H\ge5$}：半衰期太短意味着价差回复极快，交易频率上升、成本侵蚀收益。这与论文铁律二一致——本文实测的 A 股短窗口反转（$L=5$）毛利夏普 $1.08$ 被成本压至 $-0.35$，同类机制在此适用。
  \item \textbf{上界 $H\le60$}：半衰期太长意味着占用保证金数月才收敛一半，资金效率过低，且期间发生结构性断裂（如政策冲击改变均衡关系）的概率上升。
\end{itemize}

\textbf{实测}：被此步拒绝 1811 次（$10.7\%$）；最终通过 75 条。\textbf{通过者的半衰期中位数仅 5.9 交易日}，紧贴下界——这说明中国商品期货中「通过协整检验」的配对，其价差大多回复极快。这一发现对规则三有实质影响：它意味着通过筛选的配对天然是高换手的，成本敏感度高，与论文铁律二形成张力。

\subsection{S6：滚动重估}

\textbf{判据}：形成期 60 交易日、交易期 21 交易日，滚动推进，两段严格不重叠。

\textbf{理由}：若配对是在包含交易期的数据上选出的，价差的平稳性会被机械地保证，构成 look-ahead。这一「形成期与交易期严格分离」的范式来自 Gatev 等(2006) 的 $6+6$ 个月滚动设计\cite{gatev2006}。

\textbf{实现要点}（本实现的关键正确性问题）：交易期的标准化必须用\emph{滚动}均值与标准差，而非形成期常数。原因很直接——形成期的价差波动无法代表交易期的波动；若用形成期 $\sigma$ 计算 $z$，交易期一旦波动放大，$|z|$ 会虚高，导致「开仓即触止损」。本文首版实现即犯此错，成交笔数虚增、结果不可用；改为「形成期 $+$ 已过去交易期」的滚动 60 日均值与标准差后恢复正常。

\textbf{实测}：形成期数 144 期，累计 16992 次配对检验，通过 75 条。

\subsection{S7：生成可交易对表与阈值交易}

\textbf{产出字段}：形成期末日期、产业链、品种 A/B、相关性、ADF $p$、对冲比率 $\beta$、半衰期 $H$、价差均值 $\mu$、价差标准差 $\sigma$。

\textbf{对冲比率}。静态版用形成期 OLS 的 $\hat\beta$；在线版用 Kalman 滤波滚动估计 $\beta_t$\cite{elliott2005}。Kalman 的状态空间形式为
\begin{equation}
\begin{aligned}
\ln P^A_t&=\alpha_t+\beta_t\ln P^B_t+\varepsilon_t, &&\varepsilon_t\sim\mathcal N(0,\sigma_\varepsilon^2),\\
\alpha_{t+1}&=\alpha_t+\eta^{\alpha}_t,\quad \beta_{t+1}=\beta_t+\eta^{\beta}_t, &&\eta\sim\mathcal N(0,\Sigma_\eta).
\end{aligned}
\end{equation}

\textbf{必须使用一步前瞻的 $\beta_t$}。这是本实现的第二个关键正确性问题：若在时刻 $t$ 先用第 $t$ 期观测更新 $\beta$，再用更新后的 $\beta$ 计算价差 $e_t$，则 $e_t$ 被当期价格拟合，残差呈现\emph{人为的}均值回归，回测会给出虚高的夏普。本文首版即因此得到夏普 $3.51$ 的荒谬结果；改为「先记录前瞻 $\beta$、再用当期观测更新」后，夏普降至 $1.09$，与成交笔数（21 笔）相符。

\textbf{阈值规则}（把连续带宽离散化为事件规则）：

\begin{table}[htbp]
\centering\footnotesize
\caption{规则三的进出场阈值}\label{tab:pairband}
\begin{tabular}{@{}>{\raggedright\arraybackslash}p{3.0cm}>{\raggedright\arraybackslash}p{5.4cm}>{\raggedright\arraybackslash}p{5.5cm}@{}}
\toprule
\textbf{条件} & \textbf{动作} & \textbf{依据} \\
\midrule
$|z_t|>2.0$ 且空仓 & 开仓（$z>2$ 空 A 多 B；$z<-2$ 反之） & s-score 范式\cite{avellaneda2010} \\
$|z_t|<0.5$ 且持仓 & 平仓（收敛了结） & 同上 \\
$|z_t|>3.5$ 且持仓 & 强制止损（不收敛防线） & \cite{leungli2015} \\
持仓 $>3H$ 且未收敛 & 平仓（时间止损） & 资金效率 \\
实际成本升至 $c'$ & 开仓阈值改为 $2.0\times(c'/c)^{1/3}$ & 带宽随成本立方根\cite{martin2011} \\
\bottomrule
\end{tabular}
\end{table}

\textbf{风险归一}。每对按风险预算归一：设 1 倍标准差的价差变动对应 \texttt{risk\_per\_pair} 的资金敞口，则日损益
\begin{equation}
\Delta\text{PnL}_t=\text{direction}\times\Delta z_t\times\text{risk\_per\_pair}.
\end{equation}
总风险预算按同时持仓上限均分（本实现：总风险 $30\%$，上限 20 对，故每对 $1.5\%$）。若不归一，多对敞口直接相加会使净值波动失控（本文首版曾得到 $-100.96\%$ 的「最大回撤」，即净值归零）。

\textbf{实测结果}：成交 21 笔，平均持仓 11.6 天，有仓日的平均同时持仓对数 1.2，年化夏普 $1.09$、年化 $4.31\%$、最大回撤 $-4.06\%$。

\textbf{必须说明的局限}：21 笔交易在 12 年样本上不足以支撑统计推断。报告中 $t=3.99$ 的 Newey--West 统计量是按 \emph{3039 个交易日}计算的，而策略在绝大多数交易日空仓，该 $t$ 值严重高估显著性；正确的检验应基于\emph{交易笔数}（$n=21$），其功效极低。因此本节的规则三结果只能读作「筛选流程可用」，不能读作「策略有效」。这也正是下一节多重检验流程要处理的问题。

\subsection{筛选拒绝率的诊断价值}

各步拒绝率的分布本身就是信息：

\begin{table}[htbp]
\centering\footnotesize
\caption{配对筛选各步拒绝率实测（16992 次检验）}\label{tab:reject}
\begin{tabular}{@{}>{\raggedright\arraybackslash}p{2.4cm}>{\raggedright\arraybackslash}p{3.4cm}>{\raggedright\arraybackslash}p{2.4cm}>{\raggedright\arraybackslash}p{6.0cm}@{}}
\toprule
\textbf{步骤} & \textbf{判据} & \textbf{拒绝次数（占比）} & \textbf{解读} \\
\midrule
S2 & $\rho\ge0.6$ & 8760（$51.6\%$） & 产业链内仍有大量品种对走势独立，先验均衡关系薄弱 \\
S3 & ADF $p<0.05$ & 4942（$29.1\%$） & 高相关不等于协整，这一步是真正的资格门槛 \\
S5 & $H\in[5,60]$ & 1811（$10.7\%$） & 通过协整者多为极快回复（中位数 5.9 日），被下界拦下一部分 \\
S4 & $0<\phi<1$ & 1（$0.0\%$） & ADF 通过后残差基本稳定 \\
\midrule
合计通过 & — & 75（$0.44\%$） & 每 227 次检验仅有 1 条可交易对 \\
\bottomrule
\end{tabular}
\end{table}

「每 227 次检验仅 1 条通过」这一比例，正好印证论文对偶七的判断：反转类型是\emph{条件性}策略，它先要在一大批候选中筛掉绝大多数，剩下的才有资格进入交易设计。这也解释了为什么反转策略的有效样本基础天然比趋势策略薄。
